\chapter{Resultados de la evaluación técnica}
\label{chap:resultados-evaluacion}

\section{Propósito y alcance de la evaluación}

Este capítulo reúne los resultados obtenidos al evaluar los modelos de segmentación contra las anotaciones de referencia de sus conjuntos de origen. Su función es responder, con valores concretos, al objetivo específico que compromete la evaluación técnica del prototipo mediante métricas de segmentación.

El alcance de lo que aquí se informa está acotado y conviene delimitarlo antes de presentar cifra alguna. Se evalúa la coincidencia entre las máscaras que produce cada modelo y las anotaciones de referencia del conjunto correspondiente. No se evalúa la exactitud clínica de las mediciones derivadas, para lo cual sería necesario un conjunto de mediciones realizadas por profesionales sobre los mismos estudios, ni la utilidad del prototipo en un flujo de trabajo real, que corresponde a la validación cualitativa prevista para la instancia siguiente.

La distinción importa porque una segmentación con buen desempeño métrico no garantiza que las mediciones derivadas de ella sean correctas, del mismo modo que una medición correcta no implica que la interpretación clínica que se haga de ella lo sea. El Capítulo 7 desarrolla esta separación en términos conceptuales; aquí se sostiene en la presentación de los resultados.

\section{Protocolo de evaluación}

La evaluación se realiza sobre una partición retenida que no interviene en el entrenamiento ni en la selección del modelo. La separación se establece a nivel de paciente y no de corte ni de serie, de modo que ninguna imagen de un paciente presente en el conjunto de entrenamiento aparezca en el de prueba. Esta condición es determinante en imagen médica: cortes contiguos de un mismo volumen guardan una similitud muy alta, y una separación por corte produciría una estimación optimista que no describe el comportamiento sobre pacientes nuevos.

{\scriptsize
\setlength{\LTleft}{0pt}
\setlength{\LTright}{0pt}
\setlength{\tabcolsep}{3pt}
\renewcommand{\arraystretch}{1.18}
\begin{longtable}{>{\raggedright\arraybackslash}p{0.20\textwidth} >{\raggedright\arraybackslash}p{0.17\textwidth} >{\raggedright\arraybackslash}p{0.20\textwidth} >{\raggedright\arraybackslash}p{0.18\textwidth} >{\raggedright\arraybackslash}p{0.18\textwidth}}
\caption{Partición de los conjuntos de datos empleados en el entrenamiento y la evaluación.}\label{tab:res-particion-datasets}\\
\hline
\textbf{Conjunto} & \textbf{Unidad de separación} & \textbf{Entrenamiento} & \textbf{Validación} & \textbf{Prueba} \\
\hline
\endfirsthead
\hline
\textbf{Conjunto} & \textbf{Unidad de separación} & \textbf{Entrenamiento} & \textbf{Validación} & \textbf{Prueba} \\
\hline
\endhead
\hline
\endfoot
SPIDER, plano sagital & Paciente & 152 pacientes, 5.271 cortes & 33 pacientes, 1.174 cortes & 33 pacientes, 1.218 cortes \\
Sudirman y Al-Kafri, plano axial & Paciente & Pendiente de recuperación & 81 casos evaluables & 102 casos evaluables \\
\hline
\end{longtable}
}

La selección de cortes del conjunto sagital se realiza mediante una distribución lineal de hasta dieciséis cortes por volumen, conservando una proporción de cortes sin estructura anotada del veinte por ciento y un máximo de cuatro cuando el volumen no contiene ninguna. El criterio evita que el modelo se entrene únicamente sobre cortes centrales, donde las estructuras son más nítidas.

La trazabilidad del entrenamiento sagital se conserva mediante identificadores criptográficos del índice de entrenamiento, del código de la arquitectura y del commit del repositorio, junto con la referencia al cuaderno que produjo el artifact.

En el caso axial, la separación a nivel de paciente está confirmada en el manifest del artifact, pero los conteos de la partición de entrenamiento no fueron publicados en él. El proyecto conserva dos particiones axiales que no coinciden: la congelada en una etapa anterior y la derivada por el cuaderno que produjo el artifact en uso. La recuperación de los conteos definitivos permanece como tarea pendiente y se declara como tal, en lugar de informar cifras de una partición que no corresponde al modelo evaluado.

\section{Resultados de segmentación}

{\scriptsize
\setlength{\LTleft}{0pt}
\setlength{\LTright}{0pt}
\setlength{\tabcolsep}{3pt}
\renewcommand{\arraystretch}{1.18}
\begin{longtable}{>{\raggedright\arraybackslash}p{0.20\textwidth} >{\raggedright\arraybackslash}p{0.27\textwidth} >{\raggedright\arraybackslash}p{0.12\textwidth} >{\raggedright\arraybackslash}p{0.16\textwidth} >{\raggedright\arraybackslash}p{0.18\textwidth}}
\caption{Coeficiente Dice sobre la partición retenida, por modelo.}\label{tab:res-dice-modelos}\\
\hline
\textbf{Modelo} & \textbf{Clases evaluadas} & \textbf{Dice macro} & \textbf{Umbral del proyecto} & \textbf{Resultado} \\
\hline
\endfirsthead
\hline
\textbf{Modelo} & \textbf{Clases evaluadas} & \textbf{Dice macro} & \textbf{Umbral del proyecto} & \textbf{Resultado} \\
\hline
\endhead
\hline
\endfoot
Sagital, SPIDER & Vértebras, canal y discos, sin fondo & 0,893 & 0,70 & Supera el umbral \\
Axial, Sudirman y Al-Kafri & Primer plano oficial, incluyendo \texttt{raw\_0} & 0,679348283374592 (IoU 0,5699727598420694) & 0,70 & No alcanza el umbral \\
Axial, excluyendo \texttt{raw\_0} & Analisis secundario de estructuras & 0,8105661875685084 (IoU 0,6915400877715925) & 0,70 & Lectura secundaria; no reemplaza el gate \\
\hline
\end{longtable}
}

El valor informado corresponde al promedio macro del coeficiente Dice, es decir al promedio de los coeficientes por clase, que otorga el mismo peso a cada estructura con independencia de su tamaño. El criterio es más exigente que un promedio ponderado por número de píxeles, bajo el cual las estructuras pequeñas quedarían enmascaradas por las grandes.

La diferencia entre las dos lecturas del modelo axial requiere explicación, porque es de una magnitud considerable y no responde a un cambio en el modelo sino en el conjunto de clases sobre el que se promedia. Según establece la Tabla~\ref{tab:cap-clases-modelo-axial}, el modelo axial distingue seis clases, de las cuales dos corresponden a fondo. Una de ellas, la que designa la región exterior al paciente en el conjunto de origen, obtiene un coeficiente Dice de 0,15 y arrastra el promedio macro cuando se la incluye entre las clases de primer plano.

El desempeño secundario excluyendo \texttt{raw\_0} es de 0,8105661875685084 para Dice y 0,6915400877715925 para IoU. Esa lectura ayuda a describir la segmentación de las estructuras de interes, pero no sustituye la metrica oficial del artifact: \texttt{qualityGatePassed=false}, con razon \texttt{dice\_macro\_foreground\_below\_threshold}.

\section{Umbral de calidad y estado del artifact axial}

El proyecto estableció un umbral de 0,70 sobre el coeficiente Dice macro de primer plano como condición para promover un artifact a la categoría de baseline técnico. El umbral se acompaña de tres condiciones adicionales no configurables: la partición retenida debe estar separada a nivel de paciente, el informe debe reportar de forma explícita el desempeño sobre la clase de fondo del conjunto de origen, y debe incluir la intersección sobre unión además del coeficiente Dice. La generación de una partición nueva requiere además una habilitación manual.

Aplicado al modelo axial, el umbral no se alcanzó en su lectura de primer plano completo. El artifact quedó registrado como candidato por debajo del umbral de calidad y no fue promovido a baseline técnico. El valor del umbral no se modificó para acomodar el resultado.

Esta decisión merece explicitarse porque describe el papel que el control cumple en el proyecto. Un umbral que se ajusta cuando un modelo no lo alcanza no ejerce ninguna función de control; su utilidad reside precisamente en registrar los casos en que no se supera. El módulo axial, en consecuencia, se documenta como componente complementario evaluado y no como baseline consolidado, y esa condición se refleja en el estado de las capacidades del Capítulo 11.

Corresponde señalar que el criterio del umbral es discutible en un aspecto concreto: promediar una clase de fondo junto con las estructuras anatómicas penaliza al modelo por una clase que ninguna medición utiliza. La revisión de ese criterio, o su sustitución por el promedio sobre estructuras anatómicas exclusivamente, es una de las tareas previstas para la instancia siguiente.

\section{Advertencia sobre la partición de prueba axial}

La partición de prueba del conjunto axial había sido evaluada previamente durante el desarrollo de una versión anterior del mismo modelo. La evaluación que aquí se informa es, por lo tanto, comparativa respecto de esa versión y no constituye una validación sobre datos enteramente no observados.

La advertencia se registra en el manifest del artifact y el cuaderno de entrenamiento incorpora una restricción en tiempo de ejecución que impide reevaluar la partición retenida sin una confirmación explícita. Se declara aquí porque omitirla presentaría el resultado como más independiente de lo que es.

\section{Verificación por contraste de las mediciones derivadas}

Las mediciones derivadas de las máscaras no cuentan con una validación formal contra mediciones realizadas por profesionales, que requeriría un conjunto anotado específicamente para ese fin y forma parte del trabajo previsto. Durante el desarrollo se aplicó, en cambio, una verificación por contraste con rangos conocidos, cuyo alcance es más limitado pero permitió detectar errores de método que las métricas de segmentación no revelan.

{\scriptsize
\setlength{\LTleft}{0pt}
\setlength{\LTright}{0pt}
\setlength{\tabcolsep}{3pt}
\renewcommand{\arraystretch}{1.18}
\begin{longtable}{>{\raggedright\arraybackslash}p{0.27\textwidth} >{\raggedright\arraybackslash}p{0.20\textwidth} >{\raggedright\arraybackslash}p{0.27\textwidth} >{\raggedright\arraybackslash}p{0.20\textwidth}}
\caption{Verificación por contraste de las mediciones sobre el estudio de referencia.}\label{tab:res-contraste-mediciones}\\
\hline
\textbf{Medición} & \textbf{Valor obtenido} & \textbf{Referencia de contraste} & \textbf{Lectura} \\
\hline
\endfirsthead
\hline
\textbf{Medición} & \textbf{Valor obtenido} & \textbf{Referencia de contraste} & \textbf{Lectura} \\
\hline
\endhead
\hline
\endfoot
Diámetro anteroposterior del canal, mediana & 13,1 mm & Rango compatible con el saco tecal lumbar & Consistente \\
Diámetro anteroposterior por nivel, de T11-T12 a L4-L5 & 13,1 · 14,2 · 13,1 · 12,0 · 9,8 · 12,0 mm & Descenso esperable con un valor menor aislado & Consistente, con un valor menor en L3-L4 \\
Suma de ángulos segmentarios de L1 a L5 & 36 grados & Lordosis lumbar habitual entre 40 y 60 grados & Del mismo orden, por debajo del rango \\
\hline
\end{longtable}
}

Este procedimiento no equivale a una validación. Un contraste con rangos poblacionales permite descartar valores manifiestamente incompatibles, pero no establece la exactitud de una medición individual. Su utilidad durante el desarrollo fue detectar tres errores de método cuya magnitud se documenta en la sección 11.5, que las métricas de segmentación no habrían señalado, dado que la máscara era correcta y el error residía en cómo se medía sobre ella.

\section{Una medición implementada y no publicada}

El sistema calcula el desplazamiento entre cuerpos vertebrales contiguos, hallazgo estructural frecuente en la evaluación lumbar. Esa medición está implementada y se decidió no exponerla al usuario. El fundamento de la decisión es metodológico y se sostiene en valores medidos.

La medición depende en su totalidad de la ubicación de un único punto, la esquina posterior del platillo vertebral, que es precisamente la región donde la segmentación presenta menor precisión. Dos procedimientos igualmente razonables para derivar ese punto, uno a partir del eje de ancho del cuerpo vertebral y otro a partir del borde de la máscara, producen resultados que difieren hasta en 17 milímetros sobre el mismo nivel: 7,30 frente a 24,06 milímetros entre la primera y la segunda vértebra lumbar, y 6,64 frente a 20,78 entre la segunda y la tercera. Con una dispersión de ese orden, cualquiera de los dos valores resulta arbitrario.

El contraste con el ángulo segmentario, que sí se publica con carácter experimental, no reside en el método sino en la posibilidad de verificación. El ángulo pudo contrastarse contra un rango poblacional conocido, según se indica en la sección anterior. Para el desplazamiento vertebral no existe un contraste equivalente sin un conjunto de estudios con el desplazamiento medido por un profesional.

La decisión se documenta porque describe un criterio que el proyecto aplica de forma general: una medición se publica cuando puede verificarse, y no cuando puede calcularse.


\section{Resultados finales de validacion tecnica por capa}
\label{sec:resultados-validacion-final}

Ademas de la evaluacion offline de modelos, el cierre tecnico del producto incluyo verificaciones de software, integracion y recorrido browser. Estas pruebas validan el comportamiento funcional del prototipo bajo condiciones controladas; no constituyen validacion clinica.

{\scriptsize
\setlength{\LTleft}{0pt}
\setlength{\LTright}{0pt}
\setlength{\tabcolsep}{3pt}
\renewcommand{\arraystretch}{1.18}
\begin{longtable}{>{\raggedright\arraybackslash}p{0.26\textwidth} >{\raggedright\arraybackslash}p{0.28\textwidth} >{\raggedright\arraybackslash}p{0.38\textwidth}}
\caption{Validacion tecnica final del prototipo.}\label{tab:res-validacion-tecnica-final}\\
\hline
\textbf{Validacion} & \textbf{Resultado} & \textbf{Interpretacion} \\
\hline
\endfirsthead
\hline
\textbf{Validacion} & \textbf{Resultado} & \textbf{Interpretacion} \\
\hline
\endhead
\hline
\endfoot
Backend Java 21, suite reconciliada & 1053 tests, 0 failures, 0 errors, 78 skipped & Suite final verde para el alcance ejecutado; los skipped corresponden a tests condicionados por Postgres/Testcontainers. No mezclar con conteos baseline previos. \\
Backend focal & 67/67 PASS & Cubre Patient API, controlador AI, seguridad, validadores y asociacion Study--Patient a nivel Backend. \\
Frontend & lint PASS, typecheck PASS, build PASS, unit tests PASS & Valida calidad estatica y construccion de la interfaz final. \\
AI focal/product regressions & PASS; 67 passed y 1 skipped en el recorte focal & Valida product orchestration, multiplanar, real baseline, Grad-CAM y subarticular en el alcance focal. \\
AI full suite & 479 passed / 40 failed & No se presenta como suite totalmente verde: los fallos son brechas ambientales preexistentes asociadas a \texttt{highdicom} y resolucion de subprocess. \\
Docker full stack & PASS & Construccion del stack completo con imagenes requeridas. \\
Browser E2E & Playwright 11/11 PASS & Validacion funcional del producto en recorrido demo/final; no equivale a validacion clinica ni a Patient entity browser-level completo. \\
Patient API smoke & 8/8 PASS & Valida create/associate/list/detail/longitudinal/report/reload/persist-after-restart a nivel API/Backend. \\
Persistencia & PASS & Resultados, asociaciones y series sobreviven reinicio del Backend. \\
Product-analysis & PASS en demo path; bloqueado en path real sin 67B2 parity & El producto no fabrica nivel anatomico cuando la paridad runtime no esta probada. \\
Longitudinal & PASS & Comparacion deterministica en Demo Mode y flujo \texttt{subjectRef}. \\
Reporte / PDF / JSON & PASS & Reporte estructurado \texttt{pfi.study-report.v1}; exportaciones PDF y JSON verificadas. \\
Chat contextual & PASS en mock LLM & Valida integracion funcional con proveedor deterministico; no despliegue productivo de Qwen. \\
Evidencia visual & PASS & Overlays/assets servidos por Backend; limitados por registro AI efimero tras reinicio. \\
DICOM invalido & 422 \texttt{UNSUPPORTED\_INPUT} PASS & El sistema falla de forma controlada ante input no procesable. \\
Secret scan & Sin credenciales reales & Hallazgos revisados como falsos positivos; no se detectaron secretos productivos. \\
\hline
\end{longtable}
}

\section{Niveles de validacion y alcance}
\label{sec:resultados-niveles-validacion}

El cierre distingue cinco niveles. La validacion de modelo offline mide desempeno contra anotaciones o particiones congeladas. La validacion tecnica de software verifica contratos, persistencia, seguridad y errores controlados. La validacion E2E recorre el producto integrado en browser y stack containerizado. La validacion profesional/usabilidad requiere que profesionales revisen el producto integrado y emitan observaciones; esa etapa no se completa en este documento. La validacion clinica requeriria un protocolo especifico con datos, comparadores y criterios clinicos; no fue realizada y no se afirma.

Esta separacion es central para interpretar las metricas. Por ejemplo, 67B2 tiene resultados offline altos, pero conserva \texttt{67B2\_RUNTIME\_PARITY\_NOT\_PROVEN}, \texttt{productEnabled=false} y \texttt{parityVerified=false}; por lo tanto no habilita inferencia productiva real. Del mismo modo, Playwright 11/11 valida el recorrido funcional del prototipo, no la exactitud clinica de sus salidas.
\section{Limitaciones de la evaluación}
\label{sec:resultados-limitaciones}

La evaluación final permite afirmar la consistencia funcional y técnica del prototipo académico, pero no habilita conclusiones clínicas ni productivas fuera del alcance declarado. Las limitaciones principales son:

\begin{enumerate}
    \item No se realizó validación clínica prospectiva ni retrospectiva con profesionales sobre una cohorte independiente.
    \item El recorrido Playwright 11/11 valida comportamiento funcional de interfaz, no desempeño clínico.
    \item El modo demo usa fixtures determinísticos y no debe interpretarse como inferencia médica real.
    \item La integración Patient está implementada en Backend/API/persistencia, pero la interfaz conserva flujos basados en \texttt{subjectRef} y no ofrece creación/selección completa de Patient en navegador.
    \item La prueba Patient E2E 8/8 corresponde a validación API-level/smoke, no a recorrido browser-level completo.
    \item El modelo axial oficial no supera el umbral de calidad definido, por lo que queda como capacidad experimental o complementaria.
    \item La mejora axial secundaria que excluye \texttt{raw\_0} es evidencia exploratoria y no reemplaza la métrica oficial de gate.
    \item El modelo 67B2 tiene métricas offline verificadas, pero no se probó paridad runtime productiva y permanece con \texttt{productEnabled=false}.
    \item P10.7 mantiene \texttt{automaticInferenceEnabled=false} por \texttt{blocked\_preprocessing\_parity}.
    \item 67C no se declara validado productivamente.
    \item Grad-CAM cuenta con pruebas AI específicas, contrato API y wiring Frontend, pero permanece experimental y no fue ejercitado explícitamente en el recorrido browser E2E final.
    \item El chat contextual se validó con mock determinístico; el uso de Qwen queda como ruta objetivo, no como disponibilidad productiva demostrada.
    \item El despliegue Cloud está limitado a recursos DEV corroborados; AI Cloud Run privado, Cloud SQL, Secret Manager, storage productivo, DEMO environment y monitoreo/alerting real no fueron desplegados.
    \item El escaneo de secretos no encontró credenciales reales, pero no reemplaza una auditoría de seguridad formal.
    \item El sistema no aprende automáticamente de correcciones profesionales; la curación de feedback queda como insumo controlado para entrenamiento offline futuro.
\end{enumerate}
\section{Conclusión del capítulo}

El módulo sagital alcanza un coeficiente Dice macro de 0,893 sobre una partición retenida separada a nivel de paciente, supera el umbral establecido y se documenta como baseline técnico del prototipo. El módulo axial registra una métrica oficial Dice 0,679348283374592 e IoU 0,5699727598420694, por debajo del umbral; la lectura secundaria que excluye \texttt{raw\_0} alcanza Dice 0,8105661875685084 e IoU 0,6915400877715925, pero no reemplaza el gate oficial.

Los resultados sostienen que la segmentación de estructuras lumbares en el plano sagital es viable con los datos y la arquitectura empleados, y que el módulo axial requiere trabajo adicional antes de considerarse consolidado. Ninguno de los dos resultados autoriza afirmaciones sobre desempeño clínico, que dependen de una validación que este trabajo no realiza y declara como fuera de su alcance.

El aporte de este capítulo no reside únicamente en las cifras. Reside también en documentar un control de calidad que se aplicó y no se superó, una partición de prueba cuya independencia es parcial, una verificación de mediciones que no equivale a validación y una medición que se decidió no publicar. Un prototipo académico se evalúa por la correspondencia entre lo que afirma y lo que puede demostrar, y esa correspondencia exige informar también lo que no funcionó como se esperaba.
